\documentclass[10pt,a4paper]{amsart}
\usepackage[margin=19mm]{geometry}
\usepackage{amsmath,amssymb,mathtools}
\usepackage[scr=rsfs,cal=euler]{mathalfa}
\usepackage{xfrac}
\usepackage[unicode,hidelinks,bookmarks=false]{hyperref}
\newcommand{\ie}{i.e.\ }
\newcommand{\cf}{cf.\ }
\newcommand{\resp}{respectively\ }
\newcommand{\Subsection}{\subsection}
\providecommand{\nobreakdash}{\nobreak\hbox{-}}
\setlength{\emergencystretch}{2em}
\multlinegap0pt
\usepackage{bm}
\usepackage{bbm}
\usepackage{dsfont}
\usepackage{xspace}
\usepackage{cancel}
\usepackage{mathtools}
\usepackage{amsthm}
\makeatletter
\@ifundefined{proposition}{\newtheorem{proposition}{Proposition}}{}
\@ifundefined{lemma}{\newtheorem{lemma}[proposition]{Lemma}}{}
\makeatother
\title[Extensions of integral orthoregular sets and icubes]{Extensions of integral orthoregular sets and icubes}
\author{M\'arton Erd\'elyi, P\'eter Maga, Gergely Z\'abr\'adi}
\date{Source: arXiv:2508.01196v1 (CC BY 4.0)}
\begin{document}
\maketitle

\noindent\emph{Source and licence.} Extensions of integral orthoregular sets and icubes. Version arXiv:2508.01196v1 (CC BY 4.0), distributed under the Creative Commons Attribution 4.0 International licence, as recorded in the arXiv licence metadata for this exact version and shown on the abstract page https://arxiv.org/abs/2508.01196v1. Full rights notice: licenses/arxiv-2508.01196-CC-BY-4.0.txt.\par\smallskip

\noindent\textit{Editorial scope.} Counted unit: the Proposition whose source label is \texttt{factorizationthm} in the article (source0.tex lines 594--600, source label \texttt{factorizationthm}), delivered together with the article's own complete original proof of it (source0.tex lines 602--652, one \texttt{proof} environment of 51 lines). No mathematics is written, repaired or re-derived: statement and proof are the article's own text line for line. Required context retained uncounted: perpprop2 (lines 538--541); Qsquare (lines 544--546); perpconverse (lines 575--577). Every definition, display and label of those lines is retained unchanged. Omitted from this package and not redistributed: 1--537, 542--543, 547--574, 578--593, 601--601, 653--1083. Nothing inside a retained excerpt is omitted or altered: every retained body is delivered exactly as the article's lines are written, so no omission receipt of any kind accompanies this package. Citations: the retained excerpts carry no citation command at all, so no bibliography entry of the article is transcribed and no reference is redistributed. Reference adaptation: none was needed and none was made; every reference inside the delivered unit resolves inside it, so no unresolved reference or citation is produced. Printed slips kept literal: no printed word, symbol, hypothesis or number of the counted statement or of its proof was altered. Layout and class: the article's own class is \texttt{amsart} with packages including inputenc, t1enc, babel, calc, latexsym, amssymb, amsmath, amsthm, comment, hyperref, enumitem, xcolor, geometry, soul; the wrapper is the lane's pinned \texttt{amsart} editorial wrapper at 10pt on A4, which restates the theorem-like environments the retained text needs and adds the article's own macro definitions that the retained text needs (none), with extra wrapper packages (bm, bbm, dsfont, xspace, cancel, mathtools). The delivered numbering is the unit's own. Assets: no retained span contains a figure environment, a graphics-inclusion command, a PDF-page inclusion, a table or a TikZ picture; no image file is copied into this package, the package declares no asset dependency and no image file is redistributed with it. Licence: version arXiv:2508.01196v1 of this article is distributed under the Creative Commons Attribution 4.0 International licence, as recorded in the arXiv licence metadata for this exact version and shown on the abstract page https://arxiv.org/abs/2508.01196v1. A full rights notice is delivered at licenses/arxiv-2508.01196-CC-BY-4.0.txt. Characters: every retained excerpt is pure ASCII, so the delivered engine, XeLaTeX with its default Latin Modern text font, reports no missing character.
\begin{proposition}\label{perpprop2}~
If $(a_1,a_2)\in K^2\times K^2$ is a
$Q$-orthobalanced basis of norm $\lambda$, then $a_2=(Ma_1)_\perp/\delta$ for some $\delta\in K$, $|\delta|^2=\Delta$.
\end{proposition}

\begin{equation}\label{Qsquare}
\alpha Q((x,y)^T)=|\alpha x+\beta y|^2+(\alpha\gamma-|\beta|^2)|y|^2=|\alpha x+\beta y|^2+\mu|y|^2.
\end{equation}

\begin{lemma}\label{perpconverse}
Let $a_1\in K^2$, $\delta\in K$, and $a_2=(Ma_1)_\perp/\delta$ be such that $Q(a_1)=\lambda$ and $|\delta|^2=\Delta$. Then $(a_1,a_2)$ is a $Q$-orthobalanced basis of norm $\lambda$ and of type $\delta$.
\end{lemma}

\begin{proposition}\label{factorizationthm}
The map $F:K^2\times K^2\to L^2$ given by
\[
F((x,y)^T,(w,z)^T):=\left(\overline z+\overline y\sqrt{\varepsilon}j,-w+\overline x\sqrt{\varepsilon}j\right)
\]
restricts to a bijecton between the sets $\mathcal B_Q(\delta,\lambda)$ and $\mathcal F_{\frac{\alpha\lambda}{\Delta}}\left(\frac{\lambda}{\Delta}(\beta+\delta\sqrt{\varepsilon}j)\right)$.
\end{proposition}

\begin{proof}
It is obvious that $F$ is a bijection from $K^{2\times 2}$ to $L^2$.

If we restrict $F$ to $\mathcal B_Q(\delta,\lambda)$, then by Proposition~\ref{perpprop2},
\[
\begin{pmatrix}
w \\ z
\end{pmatrix} = \frac{1}{\delta}\left(M\begin{pmatrix}
x \\ y
\end{pmatrix}\right)_{\perp} = \frac{1}{\delta} \begin{pmatrix}
-\beta \overline{x} - \gamma \overline{y}\\ \alpha \overline{x} + \overline{\beta} \overline{y}
\end{pmatrix}.
\]


Setting
\[
u:=\overline z+\overline y\sqrt{\varepsilon}j,\qquad v:=-w+\overline x\sqrt{\varepsilon}j,
\]
by (\ref{Qsquare}),
\begin{equation}\label{u2eq}
|u|^2=|\overline z+\overline y\sqrt{\varepsilon}j|^2=\left|\frac{\alpha x+\beta y}{\overline\delta}\right|^2+\varepsilon |\overline y|^2=\frac1\Delta\left(|\alpha x+\beta y|^2+\Delta\varepsilon|y|^2\right)=\frac{\alpha\lambda}{\Delta}.
\end{equation}
On the other hand, using that for $s\in K$, we have $sj=j\overline s$, we see that
\begin{equation}\label{uveq}
\begin{split}
uv&=(\overline z+\overline y\sqrt{\varepsilon}j)(-w+\overline x\sqrt{\varepsilon}j)=-(\overline zw+\varepsilon\overline yx)+\overline{(zx-yw)}\sqrt{\varepsilon}j\\ 
&=\frac1\Delta\cdot\left(((\alpha x+\beta y)(\beta\overline x+\gamma\overline y)-\Delta\varepsilon\overline y x)+((\alpha x+\beta y)\overline x+\overline y(\overline \beta x+\gamma y))\delta\sqrt{\varepsilon}j\right) \\ 
&=\frac1\Delta(\beta Q(a_1) +\delta Q(a_1)\sqrt{\varepsilon}j)=\frac\lambda\Delta(\beta+\delta\sqrt{\varepsilon}j).
\end{split}
\end{equation}
Then \eqref{u2eq} and \eqref{uveq} together show that $\mathcal B_Q(\delta,\lambda)$ is indeed mapped to $\mathcal F_{\frac{\alpha\lambda}{\Delta}}\left(\frac{\lambda}{\Delta}(\beta+\delta\sqrt{\varepsilon}j)\right)$.

For the converse, assume that $(u,v)\in \mathcal{F}_{\frac{\alpha\lambda}{\Delta}}\left(\frac{\lambda}{\Delta}(\beta+\delta\sqrt{\varepsilon}j)\right)$, i.e.
\[
|u|^2=\frac{\alpha\lambda}{\Delta},\qquad uv=\frac{\lambda}{\Delta}(\beta+\delta\sqrt{\varepsilon}j),
\]
and we give explicitly $F^{-1}(u,v)\in\mathcal B_Q(\delta,\lambda)$. Writing $u=r+s\sqrt{\varepsilon} j\in L$ with $r,s\in K$, we first choose
\begin{equation}\label{eq:choice-for-u}
z:= \overline{r},\qquad y:=\overline{s},
\end{equation}
and complete this data with
\[
x:=\frac{\overline{\delta}\overline{z}-\beta y}{\alpha},\qquad w:= - \frac{\beta\overline{x}+\overline{\gamma y}}{\delta}.
\]
Applying \eqref{Qsquare}, we see that
\[\alpha Q((x,y)^T)=|(\overline\delta\overline z-\beta y)+\beta y|^2+\Delta\varepsilon|y|^2=\Delta |u|^2=\alpha\lambda,\]
hence $Q((x,y)^T)=\lambda$, and then Lemma~\ref{perpconverse} shows that $((x,y)^T,(w,z)^T)\in\mathcal B_Q(\delta,\lambda)$.

To complete the proof of this direction, we have to verify that $F((x,y)^T,(w,z)^T)=(u,v)$. As for the $u$-part, this is clear from \eqref{eq:choice-for-u} and the definition of $F$. As for the $v$-part, we have to check that $v':=-w+\overline{x}\sqrt{\varepsilon}j$ is the same as the originally given $v$. This indeed holds, since $uv=\frac{\lambda}{\Delta}(\beta+\delta\sqrt{\varepsilon}j)$ (by assumption), $uv'=\frac{\lambda}{\Delta}(\beta+\delta\sqrt{\varepsilon}j)$ (by the already proven direction), and $L$ is a division algebra.
\end{proof}

\end{document}
